% Part: incompleteness
% Chapter: introduction
% Section: historical-background

\documentclass[../../../include/open-logic-section]{subfiles}

\begin{document}

\olfileid{inc}{int}{bgr}

\olsection{خلفية تاريخية}

سنناقش في هذا القسم بإيجاز تطورات تاريخية تساعد على وضع مبرهنات عدم
الاكتمال في سياقها. وسنقدم، على وجه الخصوص، لمحة شديدة الإيجاز عن
تاريخ المنطق الرياضي، ثم نقول كلمات قليلة في تاريخ أسس الرياضيات.

\begin{digress}
  عبارة «المنطق الرياضي» ملتبسة. فيمكن فهم كلمة «رياضي» وصفًا لموضوع
  البحث، كما في «منطق الرياضيات»، أي مبادئ الاستدلال الرياضي؛ أو
  وصفًا للمناهج، كما في «رياضيات المنطق»، أي الدراسة الرياضية لمبادئ
  الاستدلال. والرواية الآتية تتناول المنطق الرياضي بالمعنيين، وكثيرًا
  ما تتناولهما معًا.
\end{digress}

بدأت دراسة المنطق، في جوهر الأمر، مع أرسطو، الذي عاش نحو 384--322
\textsc{bce}. وتشتمل كتبه \emph{المقولات} و\emph{التحليلات الأولى}
و\emph{التحليلات الثانية} على دراسات منهجية لمبادئ الاستدلال العلمي،
ومنها دراسة شاملة ومنهجية للقياس.

وساد منطق أرسطو الفلسفة المدرسية طوال العصور الوسطى؛ بل ظل كانط في
أواخر القرن الثامن عشر يرى أن منطق أرسطو كامل لا يحتاج إلى مراجعة.
لكن نظرية القياس أضيق كثيرًا من أن تمثل أكثر من أشد جوانب الاستدلال
الرياضي سطحية. وقبل ذلك بقرن تصور لايبنتس، معاصر نيوتن، «حسابًا» كاملًا
للاستدلال المنطقي، وخطا خطوات أولية في تصميمه، فقدم في جوهر الأمر
صيغة من منطق القضايا.

وكان القرن التاسع عشر منعطفًا حاسمًا للمنطق. ففي سنة 1854 كتب جورج
بول \emph{قوانين الفكر}، وهو دراسة جبرية شاملة لمنطق القضايا لا تبتعد
كثيرًا عن العروض الحديثة. وفي سنة 1879 نشر غوتلوب فريغه كتابه
\emph{الكتابة المفهومية} (Begriffsschrift)، فوسع منطق القضايا
بالكمّات والعلاقات، واشتمل بذلك على منطق الرتبة الأولى. بل إن نظم
فريغه المنطقية اشتملت أيضًا على منطق الرتبة العليا وعلى أكثر منه. وفي
\emph{القوانين الأساسية للحساب} سعى فريغه إلى بيان أن الحساب كله
يمكن اشتقاقه في كتابته المفهومية من افتراضات منطقية محضة. لكن اتضح،
للأسف، أن هذه الافتراضات غير متسقة، كما بين راسل سنة 1902. وإذا
نحينا البديهية غير المتسقة، فقد اخترع فريغه المنطق الحديث بمفرده إلى
حد بعيد، وهو إنجاز مذهل. وطور مفكرون ذوو نزعة جبرية بعد بول، منهم
بيرس وشرودر، منطق الكم مستقلين أيضًا.

لننتقل الآن إلى التطورات في أسس الرياضيات. وبالطبع ثمة تفاعل كبير
مع التطورات السابقة، لأن المنطق يؤدي دورًا مهمًا في الرياضيات. فمثلًا
طور فريغه منطقه بغرض صريح هو بيان إمكان تأسيس الرياضيات كلها على
إطاره المنطقي وحده؛ وأراد، على وجه الخصوص، بيان أن الرياضيات تتألف من
حقائق \emph{تحليلية} قبلية، لا حقائق \emph{تركيبية} قبلية كما رأى
كانط.

ويرى كثيرون أن مولد الرياضيات بمعناها الدقيق وقع عند اليونان. فقد
كانت \emph{عناصر} إقليدس، المكتوبة نحو سنة 300 قبل الميلاد، نموذجًا
ناضجًا للرياضيات اليونانية بتشديدها على الصرامة والدقة. ولا تزال
تعريفات \emph{العناصر} وبراهينها باقية على حالها تقريبًا في كتب
الهندسة المدرسية اليوم (بقدر ما لا تزال الهندسة تدرّس في المدارس).
وقد عُد هذا النموذج من الاستدلال الرياضي مثالًا للحجاج الصارم، لا في
الرياضيات وحدها بل في فروع من الفلسفة أيضًا. (بل عرض سبينوزا حججًا
أخلاقية ودينية على الطريقة الإقليدية، وهو أمر غريب المنظر!)

ابتكر نيوتن ولايبنتس حساب التفاضل والتكامل في القرن السابع عشر.
(واحتدم نزاع على السبق قرونًا، لكن معظم الباحثين اليوم يرون أن
التطورين كانا مستقلين في جملتهما.) ويشتمل هذا الحساب، مثلًا، على
استدلال في مجاميع غير متناهية لكميات متناهية في الصغر؛ وقد غذت هذه
السمات نقد الأسقف بركلي، الذي رأى أن الإيمان بالله ليس أقل عقلانية
من رياضيات عصره. واستعملت طرائق حساب التفاضل والتكامل على نطاق واسع
في القرن الثامن عشر، ومن ذلك استعمال ليونهارت أويلر حسابات تنطوي على
مجاميع غير متناهية، بنتائج باهرة.

حاول الرياضيون في القرن التاسع عشر الرد على نقد بركلي بإقامة حساب
التفاضل والتكامل على أساس أرسخ. وأفضت جهود كوشي وفايرشتراس وبولتسانو
وغيرهم إلى تعريفاتنا المعاصرة للنهايات والاستمرار والتفاضل والتكامل
بدلالة «إبسيلونات ودلتات»، أي من غير أي إحالة إلى المتناهيات في
الصغر. ثم حاول الرياضيون في أواخر القرن المضي أبعد، وشرح جميع جوانب
الحساب، ومنها الأعداد الحقيقية نفسها، بدلالة الأعداد الطبيعية.
(كرونيكر: «خلق الله الأعداد الصحيحة، وما سواها من صنع الإنسان».) وفي
سنة 1872 كتب ديديكند «الاستمرار والأعداد غير النسبية»، فبين كيف يمكن
«بناء» الأعداد الحقيقية بوصفها مجموعات من الأعداد النسبية (التي يمكن،
كما تعلم، النظر إليها بوصفها أزواجًا من الأعداد الطبيعية)؛ وفي سنة
1888 كتب «ما الأعداد وما شأنها؟»، ساعيًا إلى شرح الأعداد الطبيعية
بعبارات «منطقية» محضة. وفي سنة 1887 كتب كرونيكر «في مفهوم العدد»،
وتحدث فيه عن تمثيل جميع الكائنات الرياضية بدلالة الأعداد الصحيحة؛ وفي
سنة 1889 قدم جوزيبي بيانو بديهيات صورية رمزية للأعداد الطبيعية.

وجلبت نهاية القرن التاسع عشر جرأة جديدة في التعامل مع اللانهاية. فقد
كانت الكائنات والبنى اللانهائية (مثل مجموعة الأعداد الطبيعية) تعالج
قبل ذلك بحذر؛ ففهمت «كثرة غير متناهية» بمعنى «أي عدد تشاء»، وفهم
«يقترب في النهاية» بمعنى «يصير قريبًا بالقدر الذي تشاء». لكن جورج
كانتور بين إمكان أخذ اللانهاية بمعناها الحرفي. وأسهم عمل كانتور
وديديكند وغيرهما في إدخال الفهم العام للرياضيات على أساس نظرية
المجموعات، وهو الفهم المقبول على نطاق واسع اليوم.

وهذا يقودنا إلى تطورات القرن العشرين في المنطق والأسس. ففي سنة 1902
اكتشف راسل المفارقة في نظام فريغه المنطقي. وفي سنة 1904 برهن تسرملو
مبدأ كانتور للترتيب الجيد باستعمال ما يسمى «بديهية الاختيار»، وأثارت
مشروعية هذه البديهية نقاشًا واسعًا. وبين 1910 و1913 ظهرت المجلدات
الثلاثة من \emph{مبادئ الرياضيات} لراسل ووايتهد، موسعة برنامج فريغه
لتأسيس الرياضيات على أسس منطقية. لكن راسل ووايتهد اضطرا للأسف إلى
تبني مبدأين بدا تسويغهما بوصفهما منطقيين محضين عسيرًا: بديهية
اللانهاية وبديهية «قابلية الاختزال». وفي العقد الأول من القرن العشرين
انتقد بوانكاريه استعمال «التعريفات غير الحملية» في الرياضيات، وبدأ
براور في العقد الثاني يقترح إعادة تأسيس الرياضيات كلها على أساس
«حدسي»، يتجنب استعمال قانون الثالث المرفوع ($!A\lor\lnot !A$).

يا لها من أيام غريبة!{} يسمى برنامج رد الرياضيات كلها إلى المنطق
اليوم «المنطقانية»، ويشيع النظر إليه بوصفه قد أخفق بسبب الصعوبات
المذكورة. ويسمى برنامج تطوير الرياضيات بدلالة إنشاءات عقلية حدسية
«الحدسية»، وينظر إليه بوصفه يفرض قيودًا مفرطة على الرياضيات اليومية.
وكان ديفيد هيلبرت، أحد أكثر الرياضيين تأثيرًا في التاريخ، عند منعطف
القرن مؤيدًا قويًا للمناهج المجردة الجديدة التي أدخلها كانتور
وديديكند: «لن يطردنا أحد من الفردوس الذي خلقه كانتور لنا». وكان في
الوقت نفسه متنبهًا إلى النقد التأسيسي لهذه المناهج الجديدة (المسماة،
على نحو طريف، «كلاسيكية» الآن). واقترح طريقة للجمع بين الحسنيين:
\begin{enumerate}
\item تمثيل المناهج الكلاسيكية ببديهيات وقواعد صورية، وتمثيل المسائل
  الرياضية ب!!{formula}s في نظام بديهي.
\item استعمال مناهج «نهائية» مأمونة لإثبات اتساق هذه النظم الاستنباطية
  الصورية.
\end{enumerate}

قطع عمل هيلبرت شوطًا بعيدًا في إنجاز الهدف الأول. ففي سنة 1899 أنجز
ذلك للهندسة في كتابه الشهير \emph{أسس الهندسة}. وفي السنوات اللاحقة
عمل هو وعدد من طلابه ومعاونيه في مجالات أخرى من الرياضيات لإنجاز ما
أنجزه للهندسة. وقدم هيلبرت نفسه نظم بديهيات للحساب والتحليل. وقدم
تسرملو بديهة لنظرية المجموعات، وسعها فرانكل وسكولم وفون نيومان
وغيرهم. وفي منتصف عشرينيات القرن العشرين كان ثمة نهجان يطالبان بلقب
بدهية الرياضيات «كلها»: \emph{مبادئ الرياضيات} لراسل ووايتهد، وما
صار يعرف بنظرية مجموعات تسرملو--فرانكل.

شرع هيلبرت سنة 1921 في مشروع بحثي يرمي إلى إثبات اتساق هذه النظم.
وأعانه فيه عدد من طلابه، ولا سيما برنايس وأكرمان، ثم غنتسن. وكانت
الفكرة الأساسية لبلوغ هذا الهدف صوغ سؤال إمكان وجود !!{derivation}
لتناقض في الرياضيات بوصفه مسألة توافقية عن متتاليات محتملة من الرموز؛
أي متتاليات محتملة من الجمل تستوفي معيار كونها !!{derivation}
صحيحًا، مثلًا، لـ$!A\land\lnot !A$ من بديهيات نظام بديهي للحساب أو
التحليل أو نظرية المجموعات. وكان برهان استحالة متتالية كهذه من
الرموز، لما كان هو نفسه برهانًا رياضيًا، قابلًا للصوغ داخل هذه النظم
البديهية. وبعبارة أخرى، توجد جملة ما $\OCon$ تنص على اتساق الحساب
مثلًا. وفوق ذلك ينبغي أن تكون هذه الجملة قابلة للبرهنة في النظم
المعنية، ولا سيما إذا لم يتطلب برهانها إلا وسائل «نهائية» شديدة
التقييد.

ويمكن ضبط الهدف الثاني، وهو أن تحسم نظم البديهيات المطورة كل مسألة
رياضية، بطريقتين. وفق الأولى نقول: لكل جملة~$!A$ في لغة نظام بديهي
للرياضيات، إما أن تكون $!A$ قابلة للبرهنة من البديهيات وإما أن تكون
$\lnot !A$ كذلك. ولو صح هذا لما وجدت جمل لا يمكن إثباتها ولا دحضها
استنادًا إلى البديهيات، ولا أسئلة لا تحسمها البديهيات. ويسمى النظام
البديهي الذي له هذه الخاصية \emph{تامًا}. وبالطبع قد يظل تحديد أي
البديلين حاصلًا مهمة عسيرة عند كل جملة بعينها. لكن ينبغي أن توجد، من
حيث المبدأ، طريقة لفعل ذلك. والواقع أن التمام كان سيستلزم وجود طريقة
كهذه في نظم البديهيات وال!!{derivation} التي تناولها هيلبرت، وإن لم يدرك
هيلبرت ذلك. أما الطريقة الثانية لفهم السؤال فهي اشتراط أقوى: وجود
طريقة آلية حسابية تحدد، لكل جملة~$!A$، ما إذا كانت قابلة للاشتقاق من
البديهيات.

برهن غودل سنة 1931 «مبرهنتي عدم الاكتمال»، فبين أن هذا البرنامج لا
يمكن أن ينجح. فلا يوجد نظام بديهي للرياضيات يكون تامًا؛ وعلى وجه
الخصوص، الجملة التي تعبر عن اتساق البديهيات جملة لا يمكن إثباتها ولا
دحضها.

ووجه ذلك ضربة قاتلة لبرنامج هيلبرت الأصلي. لكنه، كما يحدث كثيرًا في
الرياضيات، فتح أيضًا سبلًا جديدة مثيرة للبحث. فإذا لم يوجد نظام صوري
واحد شامل للرياضيات كلها، كان من المعقول تطوير نظم أضيق نطاقًا وبحث
ما يمكن إثباته فيها. ومن المعقول أيضًا تطوير مناهج برهان أقل تقييدًا
لإثبات اتساق هذه النظم، وإيجاد طرائق لقياس صعوبة إثبات اتساقها. ولما
كان غودل قد بين أن لكل نظام صوري تقريبًا مسائل لا يستطيع حسمها، كان
من المعقول البحث عن مسائل «مهمة» لا يستطيع نظام صوري معطى حسمها،
وتحديد مقدار القوة التي يحتاج إليها النظام كي يحسمها. ولا يزال
المناطقة إلى اليوم يتابعون هذه الأسئلة في فرع رياضي جديد هو نظرية
البراهين.

\end{document}
